\section{Concentration of the error factor and its phase}
\label{sec:phase}

For a nonzero complex number, its phase is its argument: the angle
it makes with the positive real axis. The factor \(L_N(z)\) in
\eqref{eq:error-factorization} has the known direction
\((-1)^Nz^{2N}\) and a positive amplitude. If \(E_N(z)\) is close
to one, the full error has nearly that direction. This information
is stronger than an absolute error bound centered at zero.

\subsection{An exponential moment after separating the source point}

Condition on \(s=S_0\) in \eqref{eq:error-random-variable} and
rescale the other coordinates by \(1-s\). The simplex Jacobian
gives density \(2N(1-s)^{2N-1}\) for \(s\); its mean is
\(1/(2N+1)\). The rescaled coordinates are independent of \(s\).
Thus
\begin{equation}
 X_N=sU+(1-s)V_N,\qquad V_N=\sum_{j=1}^NW_jx_{j,N},
 \label{eq:beta-split}
\end{equation}
where \(U,s,V_N\) are independent and the grouped weights \(W_j\)
have Dirichlet parameters \((2,\ldots,2)\). Explicitly, for
nonnegative integers \(k_j\), with \(k=\sum_jk_j\),
\begin{equation}
 \mathbb E\prod_jW_j^{k_j}=\frac{\prod_j(k_j+1)!}{(2N)_k}.
 \label{eq:dirichlet-moments}
\end{equation}
This follows by expanding each power of the sum of its two simplex
coordinates and using \eqref{eq:simplex-moments}.

Let \(A_k\) be the coefficient of \(v^k\) in
\(\prod_j(1-vx_{j,N})^{-2}\). Formula~\eqref{eq:dirichlet-moments}
gives \(\mathbb E V_N^k/k!=A_k/(2N)_k\). Symmetry of the nodes
makes all odd \(A_k\) zero and all even \(A_k\) nonnegative:
the product pairs as \(\prod_{x_j>0}(1-v^2x_j^2)^{-2}\).
If \(\ell\ge\max_j|x_{j,N}|\) and \(|t|\ell<2N\), then
\begin{equation}
 \mathbb E\e^{tV_N}
 \le\prod_j(1-tx_{j,N}/(2N))^{-2}
 \le\exp\left\{\frac{t^2T_2}
 {(2N)^2(1-t^2\ell^2/(2N)^2)}\right\}.
 \label{eq:node-exponential}
\end{equation}
The first inequality replaces \((2N)_k\) by \((2N)^k\) in
a nonnegative even series. For the second, pair the nodes and use
\(-\log(1-v)\le v/(1-v)\), \(0\le v<1\). Only bounded node
variables occur in these reciprocal series.

The source point \(U\) is unbounded. Its exponential moment is
controlled directly by \eqref{eq:weighted-cosh}. Symmetry and
concavity of the power \(t/4\) give, for \(0\le t\le4\),
\[
 \mathbb E\e^{tU}\le(\mathbb E\e^{4U})^{t/4}
 =\bigl(\mathbb E\cosh(4U)\bigr)^{t/4}\le(2N+12)^{t/4}.
\]
Convexity in \eqref{eq:beta-split} gives
\(\e^{tX_N}\le s\e^{tU}+(1-s)\e^{tV_N}\). Consequently,
for \(0\le Y\le1\) and \(Y\ell<N\),
\begin{equation}
 \mathbb E\e^{2YX_N}\le
 \frac{(2N+12)^{Y/2}}{2N+1}
 +\frac{2N}{2N+1}
 \exp\left\{\frac{Y^2T_2}{N^2(1-Y^2\ell^2/N^2)}\right\}.
 \label{eq:full-exponential}
\end{equation}
The small mean \(\mathbb E s=1/(2N+1)\) is why the unbounded
source contribution remains controlled as degree grows.

\subsection{Explicit constants depending only on degree}

For \(N\ge2\), define
\begin{equation}
 \ell_N=2\sqrt{B_{N-2}}=\tfrac12\arcosh(2N+8),\qquad
 T_N^*=2(N-1)B_{N-2};\qquad \ell_1=T_1^*=0.
 \label{eq:node-allowances}
\end{equation}
The largest absolute row sum of \(J_N\) is at most \(\ell_N\),
so \(\max_j|x_{j,N}|\le\ell_N\), \(T_2\le T_N^*\), and
\(T_4\le\ell_N^2T_N^*\). Also \(\ell_N<N\): for \(N\ge2\),
\(\cosh(2N)\ge1+2N^2+2N^4/3>2N+8\); the case \(N=1\)
is immediate. Hence every denominator below is strictly positive
on \(0\le Y\le1\). Set
\begin{equation}
 \mathcal A_N(Y)=\frac{(2N+12)^{Y/2}}{2N+1}
 +\frac{2N}{2N+1}
 \exp\left\{\frac{Y^2T_N^*}{N^2(1-Y^2\ell_N^2/N^2)}\right\}.
 \label{eq:A-budget}
\end{equation}
The envelope bounds \(m_2^{[N]}\le B_N\) and
\(m_4^{[N]}\le3B_{N+1}^2\). Thus
\begin{equation}
 \begin{aligned}
 \mathcal V_N&=\frac{2(B_N+T_N^*)}{(2N+1)_2},\\
 \mathcal Q_N&=\frac{72B_{N+1}^2+24B_NT_N^*+12(T_N^*)^2
                         +12\ell_N^2T_N^*}{(2N+1)_4}
 \end{aligned}
 \label{eq:moment-budgets}
\end{equation}
bound \(\mathbb E X_N^2\) and \(\mathbb E X_N^4\), respectively.
All these are explicit elementary functions of the integer \(N\).
They require no uncomputed recurrence coefficients. At \(N=1\),
\(\mathcal A_1(Y)=14^{Y/2}/3+2/3\),
\(\mathcal V_1=B_1/6\), and \(\mathcal Q_1=B_2^2/5\).

\begin{theorem}[A quantitative error-phase bound]
\label{thm:phase}
For every \(N\ge1\), \(0\le Y\le1\), and \(|\Im z|\le Y\),
\begin{equation}
 |E_N(z)-1|\le\epsilon_N(|z|,Y),\qquad
 \epsilon_N(r,Y)=\frac{r^2}{2}
                    \sqrt{2\mathcal A_N(Y)\mathcal Q_N}.
 \label{eq:phase-bound}
\end{equation}
There is a constant independent of \(Y\in[0,1]\) such that
\(\epsilon_N(|z|,Y)=O(|z|^2\log^2(N+1)/N)\). In particular,
if \(r_N\log(N+1)/\sqrt N\to0\), then \(E_N\to1\) uniformly
on \(\{|z|\le r_N,\ |\Im z|\le1\}\).
\end{theorem}

\begin{proof}
Taylor's integral remainder gives, for real \(v\),
\[
 |\e^{\ii zv}-1-\ii zv|\le\tfrac12|z|^2v^2\e^{Y|v|}.
\]
The linear term averages to zero by symmetry. Cauchy--Schwarz gives
\[
 |E_N(z)-1|\le\tfrac12|z|^2
  (\mathbb E X_N^4)^{1/2}(\mathbb E\e^{2Y|X_N|})^{1/2}.
\]
Since \(\e^{2Y|v|}\le2\cosh(2Yv)\), symmetry and
\eqref{eq:full-exponential} bound the last expectation by
\(2\mathcal A_N(Y)\). This proves \eqref{eq:phase-bound}.

The inequality \(\arcosh v\le\log(2v)\) gives
\(B_{N+1}=O(\log^2(N+1))\), \(\ell_N=O(\log(N+1))\),
and \(T_N^*=O(N\log^2(N+1))\). It follows that
\(\mathcal Q_N=O(\log^4(N+1)/N^2)\). In \eqref{eq:A-budget}
the first term is uniformly \(O(N^{-1/2})\); the exponent in
the second is uniformly \(O(\log^2(N+1)/N)\). Hence
\(\mathcal A_N(Y)\to1\) uniformly for \(0\le Y\le1\), and the
claimed rate follows. The stated growing-region consequence is direct.
\end{proof}

\begin{corollary}[Real-axis signs and adjacent brackets]
\label{cor:brackets}
For real \(x\),
\begin{equation}
 1-\frac{x^2\sigma_N^2}{2}\le E_N(x)\le1,\qquad
 |L_N(x)|(1-x^2\mathcal V_N/2)
 \le(-1)^N(F(x)-G_N(x))\le|L_N(x)|.
 \label{eq:real-sign}
\end{equation}
Define the explicit sign radius
\begin{equation}
 R_N^{\rm sign}=\sqrt{\frac{(2N+1)(2N+2)}{B_N+T_N^*}}.
 \label{eq:sign-radius}
\end{equation}
For \(0<|x|<\min(R_N^{\rm sign},R_{N+1}^{\rm sign})\), the
even member of \(G_N,G_{N+1}\) lies strictly below \(F\), and the
odd member strictly above it. All three equal one at zero.
\end{corollary}

\begin{proof}
Symmetry gives \(E_N(x)=\mathbb E\cos(xX_N)\). Apply
\(1-v^2/2\le\cos v\le1\), then
\(\sigma_N^2\le\mathcal V_N\) and
\eqref{eq:error-factorization}. The open radius is exactly
\(\sqrt{2/\mathcal V_N}\).
\end{proof}

Where \(\epsilon_N<1\), the disk \(|E_N-1|\le\epsilon_N\)
also gives \(\Re E_N>0\) and
\(|\arg E_N|\le\arcsin\epsilon_N\), by the tangents from the
origin to that disk. In particular \(F-G_N\) has no zero there
except its order-\(2N\) zero at zero. This conclusion concerns
zeros of the \emph{error}; it does not locate zeros of \(F\).
